\section{Summary and Conclusion}\label{sec:conclusion}
While the behavior of CNNs to adversarial data has generated some intrigue in computer vision since the work of \cite{szegedy2013intriguing}, its effects on deeper networks have not been explored well. We observe that adversarial perturbations for hidden layer activations generalize across different samples and we leverage this observation to devise an efficient regularization approach that could be used to train very deep architectures.  Through our experiments and analysis we make the following observations: (1) Contrary to recent methods which are inconclusive about the role of perturbing intermediate layers of a DNN in adversarial training, we have shown that for very deep networks, they play a significant role in providing a strong regularization (2) The aggregated adversarial effect of perturbing intermediate layer activations is much stronger than perturbing only the input (3) Significant improvement in classification accuracy entails capturing more variations in the data distribution while adversarial robustness can be improved by suppressing the unnecessary variations learned by the network \ref{subsec:toy-eg}. By providing an efficient adversarial training approach that could be used to regularize very deep models, we hope that this can inspire more robust network designs in the future. 

% In this work, we have attempted to present an efficient framework which proposes a simple and novel way of generating adversarial perturbations that could be used to train very deep architectures without additional training overhead. An important consequence of our study is that effect of adversarially perturbing intermediate layers.  We hope that this can open up potential theoretical directions studying adversarial training from the point of view of perturbing feature space as opposed to using it as a data augmentation technique. 